% WHAT A RUNAWAY SHOWS OF A DEFINITION.
%
% `Runaway definition?' writes the parameter text the definition had
% matched, then `->', then as much of the body as it had read. Each
% parameter of that TEXT sets the character the rest of the list is written
% with, as it does anywhere a list is shown -- so a definition whose
% parameters are written with a caret pair shows them that way in the body
% too:
%
%   Runaway definition?
%   ^^C1^^C2->[^^C1|^^C2]
%
% This engine wrote `#' throughout. trip line 364 writes
% \def\a^^C1{\d#1\d\l{#2}\l#1\par\a^^@^^@a and runs away in it.
%
% The words under `You can't use a prefix with ...' name the three prefixes
% TeX82 has and not the \protected eTeX adds.
\catcode`\{=1 \catcode`\}=2 \catcode`\#=6 \catcode`\^=7
\catcode`\^^C=6
\tracingonline=1
\outer\def\bad{}
\message{[A a runaway definition whose parameter character is a caret pair]}
\def\m^^C1^^C2{[^^C1|^^C2]\bad}
\message{[B one with the usual #]}
\def\n#1#2{[#1|#2]\bad}
\message{[C a prefix in front of something it cannot prefix]}
\global\relax\global A
\message{[done]}
\end
